\section{Introduction}
\label{sec:intro}

\IEEEPARstart{W}{aste} can be recycled, composted or safely disposed of only if its materials are kept apart. India's Solid Waste Management Rules, 2016 require every waste generator to segregate waste at source into bio-degradable, non-bio-degradable and domestic hazardous streams (Rule~4), with green, white and black bins for bio-degradable, recyclable and other waste (Rule~15(h))~\cite{moefcc2016solid}. A camera that names each item's material could support this segregation at a smart bin, sorting line or collection point, and artificial intelligence is already an established tool in solid waste management~\cite{abdallah2020artificial,fang2023artificial}.

In the literature we reviewed, however, most vision work on waste sorting solves an easier problem: classifying one photograph of one object, often on a plain background, as with TrashNet~\cite{yang2016classification} and its classifiers~\cite{kaya2023classification}. A heap, a conveyor or a street holds many items, some touching or partly hidden, and each must be located as well as named: an object-detection task. Majchrowska et al.~\cite{majchrowska2022deep}, the starting point of this work, localize litter with one network and classify each detected crop with another. Multi-object waste detection is thus established, but it is a smaller strand than single-object classification, and several detection studies still rely on single-object studio images: the task and scene gap (G1).

Detectors are also sensitive to training hyperparameters such as the learning rate, momentum, weight decay and augmentation strength; the TrashNet authors blamed the 22\% test accuracy of their convolutional neural network (CNN) on unsuitable hyperparameters~\cite{yang2016classification}. Metaheuristics have tuned CNNs, for example particle swarm optimization (PSO)~\cite{kennedy1995particle} for a mammography classifier~\cite{aguerchi2024cnn} and a beluga whale optimizer for an InceptionV3 waste classifier~\cite{sayed2024intelligent}, but the evidence is weaker than the reported gains suggest. In the studies we reviewed, metaheuristic tuning in waste vision targeted classifiers, and no waste detector had its training hyperparameters tuned by a swarm optimizer (G2). Tuned models were compared with other metaheuristics or with nothing, not with random search~\cite{bergstra2012random} at the same number of evaluations, and test results were single numbers without confidence intervals (G3), so a gain cannot be credited to the optimizer's guidance rather than to the search effort itself.

The data are a second weakness. Reviews find many waste-sorting vision studies set in simplified environments with artificially collected data~\cite{lu2022computer} that differ from real waste~\cite{fang2023artificial}, and accuracy depends heavily on the dataset: the same YOLOv5x, trained and tested on each dataset separately, scored 84.9 AP$_{50}$ on PlastOPol but 63.3 on TACO~\cite{cordova2022litter}. A tuning gain on studio images may thus say little about street photographs. For India, where segregation at source is mandatory, the closest resource, the Dense Waste Segmentation Dataset (DWSD)~\cite{ali2025dwsd}, reports segmentation pixel accuracy only, and we found no Indian real-world waste test set with published object-detection baselines (G4).

SwarmSort tests the tuning question under strict conditions. It trains YOLOv8n~\cite{jocher2023ultralytics}, a nano detector small enough for a smart bin or a phone, and compares three ways (arms) of choosing six of its training hyperparameters: the library defaults (Arm~A), the best of 32 random samples (Arm~B) and a hand-written PSO with 8 particles over 4 rounds (Arm~C). Both searches spend 32 short proxy trainings. Each chosen setting is trained once for 100 epochs and tested once on a held-out split, with paired-bootstrap 95\% confidence intervals on every difference. \genone used studio-style images in six classes; \gentwo added real-world photographs, a seventh class (OTHER) and a test-only Indian set built from DWSD. We address three research questions:
\begin{enumerate}
\item[RQ1:] At an equal budget of 32 proxy trainings, does PSO find YOLOv8n training hyperparameters that beat random search and the default settings on a held-out test set? (G2, G3)
\item[RQ2:] Does a short proxy training (12 epochs on 40\% of the data) rank hyperparameter settings the same way longer training does? (G3)
\item[RQ3:] Does the answer to RQ1 change when the data move from studio images to real-world and Indian scenes, and how large is the studio-to-real-world drop? (G1, G4)
\end{enumerate}

Our contributions are as follows.
\begin{itemize}
\item \textit{A lightweight multi-object waste detector} (G1): YOLOv8n trained on multi-object studio and real-world scenes, with six classes in \genone and seven in \gentwo (Sections~\ref{subsec:datasets} and~\ref{subsec:detector}).
\item \textit{A swarm-tuned waste detector} (G2): a hand-written PSO searches the learning-rate schedule, momentum, weight decay and two augmentation settings (Sections~\ref{subsec:space} and~\ref{subsec:pso}); to our knowledge, within the literature we reviewed, no earlier waste detector had its training hyperparameters tuned by a swarm optimizer.
\item \textit{An equal-budget, three-arm protocol} (G3): a proxy-validity check with a pass threshold fixed in advance, passed with Spearman $\rho = 0.976$ ($n = 8$, \genone only), a test split used once and paired-bootstrap 95\% confidence intervals on every difference (Sections~\ref{subsec:design}, \ref{subsec:proxy} and~\ref{subsec:eval}).
\item \textit{A two-generation evaluation across domains} (G4): studio, real-world and Indian results reported separately, including detection results on DWSD, which, to our knowledge, within the literature we reviewed, had no published detection baselines. Our DWSD boxes come from mask regions that merge touching items, and the class mapping was inferred from the masks, so these results also measure a labeling mismatch. On the same test images, real-world training data raised real-world \mapfifty from 0.038--0.096 (\genone) to 0.338--0.353 (\gentwo), at a cost of 0.074--0.105 on studio images, all significant (Section~\ref{subsec:res-domain}).
\item \textit{An honest answer, with null results reported} (O5): PSO never meaningfully beat random search; its only significant edge, $+0.003$ \mapfifty on the Indian set, lies on scores near zero. Tuning gains over the defaults were absent in \genone, small in \gentwo ($+0.017$ and $+0.020$ on the mixed test set) and did not reach real-world or Indian photographs. We explain these results (Section~\ref{sec:discussion}) and release the code, search logs, result files and trained models.
\end{itemize}

Section~\ref{sec:related} reviews related work and derives gaps G1--G4, Section~\ref{sec:methods} describes the methods, and Section~\ref{sec:results} answers RQ2, RQ1 and RQ3, in that order. Section~\ref{sec:discussion} discusses the results, their implications and the limitations, Section~\ref{sec:future} gives a prioritized roadmap, and Section~\ref{sec:conclusion} concludes.
